El apartado ``Requisitos de las acreditaciones internacionales EUR-ACE y ABET'' de la Normativa de TFT de la ETSIT-UPM establece que ``\textit{La memoria del TFT del GITST, GIB y MUIT, y en general la de aquellas titulaciones que hayan obtenido o para las que se desee solicitar una acreditación internacional EUR-ACE o ABET, … debe incluir un presupuesto económico}''.

\vspace{0.3cm}
\noindent A continuación se detalla el presupuesto económico del proyecto. El coste de mano de obra se ha calculado a partir de los 12 créditos ECTS del TFG, a razón de 25 horas por crédito (mínimo estipulado: 300 horas), más 20 horas adicionales correspondientes a reuniones con el tutor, revisiones y preparación de la defensa, totalizando 320 horas. El precio por hora de 25\,\euro\ corresponde al coste estándar de un ingeniero junior en el sector tecnológico en España. Los recursos materiales incluyen el ordenador personal del alumno (PC1, 1.500\,\euro) y el equipo de laboratorio de la ETSIT-UPM empleado para las pruebas de validación (PC2: Intel Core i9-10900X, 128\,GB DDR4, NVMe 932\,GB, precio estimado de mercado en el momento de adquisición: 2.200\,\euro). Ambos equipos se amortizan en 6 años con 8 meses de uso efectivo durante el desarrollo del proyecto.

\begin{table}[H]
    \centering
    \resizebox{\textwidth}{!}{
\begin{tabular}{lllllc}
\cline{4-6}

\multicolumn{3}{l|}{\textbf{COSTE DE MANO DE OBRA (coste directo)}} &
\multicolumn{1}{c|}{\textbf{Horas}} & \multicolumn{1}{c|}{\textbf{Precio/hora}} &
\multicolumn{1}{c|}{\textbf{Total}} \\ \cline{4-6} & &

\multicolumn{1}{l|}{} &
\multicolumn{1}{c|}{320} &
\multicolumn{1}{c|}{25 \euro} &
\multicolumn{1}{c|}{\textbf{8.000,00 \euro}} \\
\cline{4-6} & & & & &

\multicolumn{1}{l}{} \\ \cline{3-6}
\multicolumn{2}{l|}{\textbf{COSTE DE RECURSOS MATERIALES (coste directo)}} &
\multicolumn{1}{c|}{\textbf{Precio de compra}} &
\multicolumn{1}{c|}{\textbf{Uso en meses}} &
\multicolumn{1}{c|}{\textbf{Amortización (en años)}} &
\multicolumn{1}{c|}{\textbf{Total}} \\
\hline

\multicolumn{2}{|l|}{PC1: Ordenador personal del alumno} &
\multicolumn{1}{c|}{1.500,00 \euro} &
\multicolumn{1}{c|}{8} &
\multicolumn{1}{c|}{6} &
\multicolumn{1}{c|}{166,67 \euro} \\
\hline

\multicolumn{2}{|l|}{PC2: Equipo de laboratorio ETSIT-UPM (i9-10900X, 128\,GB DDR4)} &
\multicolumn{1}{c|}{2.200,00 \euro} &
\multicolumn{1}{c|}{8} &
\multicolumn{1}{c|}{6} &
\multicolumn{1}{c|}{244,44 \euro} \\
\hline

& & & & &
\multicolumn{1}{l}{} \\
\hline

\multicolumn{5}{|l|}{\textbf{COSTE TOTAL DE RECURSOS MATERIALES}} &
\multicolumn{1}{c|}{\textbf{411,11 \euro}} \\
\hline

& & & & &
\multicolumn{1}{l}{} \\
\hline

\multicolumn{1}{|l|}{\textbf{GASTOS GENERALES (costes indirectos)}} &
\multicolumn{1}{c|}{15 \%} &
\multicolumn{3}{c|}{sobre CD} &
\multicolumn{1}{c|}{\textbf{1.261,67 \euro}} \\
\hline

\multicolumn{1}{|l|}{\textbf{BENEFICIO INDUSTRIAL}} &
\multicolumn{1}{c|}{6 \%} &
\multicolumn{3}{c|}{sobre CD + CI} &
\multicolumn{1}{c|}{\textbf{580,37 \euro}} \\
\hline

& & & & &
\multicolumn{1}{l}{} \\

\multicolumn{6}{l}{\textbf{MATERIAL FUNGIBLE}} \\
\hline

\multicolumn{5}{|l|}{Impresión de la memoria} &
\multicolumn{1}{c|}{\textbf{80,00 \euro}} \\
\hline

\multicolumn{5}{|l|}{Encuadernación} &
\multicolumn{1}{c|}{\textbf{30,00 \euro}} \\
\hline

& & & & &
\multicolumn{1}{l}{} \\
\hline

\multicolumn{5}{|l|}{\textbf{SUBTOTAL PRESUPUESTO}} &
\multicolumn{1}{c|}{\textbf{10.363,15 \euro}} \\
\hline

\multicolumn{4}{|l|}{\textbf{IVA APLICABLE}} &
\multicolumn{1}{c|}{21 \%} &
\multicolumn{1}{c|}{\textbf{2.176,26 \euro}} \\
\hline

& & & & &
\multicolumn{1}{l}{} \\
\hline

\multicolumn{5}{|l|}{\textbf{TOTAL PRESUPUESTO}} &
\multicolumn{1}{c|}{\textbf{12.539,41 \euro}} \\
\hline

\end{tabular}}
    \caption{Presupuesto económico del proyecto.}
    \label{presupuesto-economico}

\end{table}
